% SCOPE: two insight-driven subsections (NOT enumerations): 2.1 RLVR/GRPO +
% Gandhi + Kim self-distill/RLRT = RL amplifies what exists; 2.2 calibration
% vs. functional metacognition (Kadavath, verbalized confidence, self-
% correction skepticism, STaR/SCoRe).
\section{Related Work}
\label{sec:related}

\subsection{RL amplifies what exists: the missing metacognitive substrate}
\label{sec:related:rlvr}

The dominant recipe for reasoning post-training is reinforcement learning
from verifiable rewards (RLVR) driven by a single terminal correctness
signal. \citet{shao2024deepseekmath} introduced GRPO, the critic-free
group-relative policy-gradient method that made outcome-reward RL practical
at scale for mathematical reasoning; DeepSeek-R1 \citep{deepseekr1} showed
that applying this recipe to a strong base model elicits long chains of
thought in which reflective-looking behavior---re-checking,
``aha''-style backtracking---\emph{emerges} without ever being rewarded as
such; SimpleRL-Zoo \citep{zeng2025simplerlzoo} replicated zero-RL across a
zoo of open base models and found that both the gains and the emergent
reflective behavior vary sharply with the base model family; and the Qwen3
report \citep{qwen3} folds the same pipeline into production post-training.
Across this entire line the objective contains one bit: was the final answer
correct. Credit is smeared over the whole trajectory, so a verification step
that actually flipped the outcome is indistinguishable, from the optimizer's
point of view, from decorative hedging that happened to co-occur with a
correct answer. Process supervision \citep{lightman2024letsverify} sharpens
credit to individual steps and outperforms outcome supervision on MATH, but
it scores steps against \emph{externally labeled} step correctness; it never
asks whether a step changed the \emph{model's own} evidence about the
answer, and so it cannot separate a step that is merely correct from one
that is epistemically useful.

Two recent lines expose, from opposite directions, why outcome-level RLVR
does not install the missing epistemic machinery. \citet{gandhi2025cognitive}
show that four cognitive behaviors---verification, backtracking, subgoal
setting, and backward chaining---predict which base models benefit from RL,
and that RL \emph{amplifies behaviors already present} in the base model or
its SFT prime: priming with the behaviors enables gains even when the primed
solutions contain wrong answers. This establishes that the substrate must be
supplied before RL can grow it, but it stops at the level of surface form:
priming installs the \emph{shape} of a habit, and amplification is
indiscriminate---nothing in the pipeline checks that an amplified
verification actually verifies. \citet{kim2026selfdistill} probe the
complementary failure: naively turning the model's own outputs into a
training signal, via self-distillation on top of RLVR, can \emph{degrade}
reasoning relative to plain GRPO, so ``the model
teaching itself'' is not automatically benign. Yet the same authors' RLRT
variant (``Rebellious Student,'' which drives exploration by \emph{reversing}
a frozen teacher's signals) \citep{kim2026rebellious} shows that a frozen
reference model's likelihoods carry genuinely useful training information
when pointed at the right target, improving over GRPO on both MATH500 and
AIME24 (numbers quoted in Section~\ref{sec:discussion}). What RLRT
does not do is localize this signal: the reference-model evidence steers
whole-trajectory exploration, never a specific reasoning act.

Read together, these results delimit our design space. Outcome-only RLVR
cannot select for functional metacognition because its credit never reaches
the metacognitive act; behavioral priming supplies forms that RL will
amplify whether or not they function; and reference-model evidence is a
double-edged signal---harmful when imitated wholesale, helpful when it
directs learning. \pmishift{} occupies the point these lines triangulate: we
prime the form (explicit \metaopen{}\dots\metaclose{} blocks via SFT), then
reward each individual block by the gold-versus-decoy evidence shift it
produces under a frozen reference model, routing that reward onto the
block's tokens only (Section~\ref{sec:method}). This is self-distillation in
the sense studied by \citet{kim2026selfdistill}, but of an evidence
\emph{contrast} (the model's own gold-versus-decoy margin) rather than of
the output distribution itself; and a matched-base twin---identical
meta-stripped data, byte-audited identical non-meta hyperparameters,
correctness-only GRPO---measures whether the installed substrate causally
pays rather than merely survives.

\subsection{Calibration is not metacognition: from knowing uncertainty to using it}
\label{sec:related:calibration}

A long line of work establishes that models \emph{know} more about their own
uncertainty than they act on. \citet{guo2017calibration} made the
miscalibration of modern networks measurable (expected calibration error,
which we adopt for RQ4), and \citet{kadavath2022language} showed that large
LMs are mostly well-calibrated when asked to judge whether their own answers
are true (P(True), P(IK))---but that knowledge is elicited by an external
probe \emph{after} generation and never routed back into the solution
process. Verbalized confidence widens the elicitation surface without
changing its role: \citet{lin2022teaching} showed models can be trained to
express calibrated uncertainty in words, and \citet{tian2023just} that
prompting RLHF-trained models yields verbalized confidences better
calibrated than their token probabilities. In all of these, calibration is
the \emph{terminal} object of study: the metric rewards a model for
accurately reporting its doubt, not for resolving it. A perfectly calibrated
model can correctly doubt every hard problem and still fail all of them.

The converse line asks whether models can \emph{use} self-assessment, and
its record is cautionary. Self-Refine \citep{madaan2023selfrefine} reported
that a model critiquing and revising its own outputs improves them across
diverse generation tasks, but \citet{huang2024llmcannot} showed that on
reasoning benchmarks, without external feedback, such intrinsic
self-correction fails to improve and can even degrade accuracy---prompting
alone does not convert latent self-knowledge into better answers.
Training-based approaches go further. STaR \citep{zelikman2022star}
bootstraps by fine-tuning on self-generated rationales filtered by answer
correctness, which reinforces whole successful traces but assigns no credit
to any epistemic act inside them. SCoRe \citep{kumar2024score} trains
multi-turn self-correction with RL and reward shaping, demonstrating that
correction \emph{behavior} is trainable---yet its reward attaches to the
outcome of the revised attempt, so credit lives at episode granularity:
correction happens between attempts, after a complete solution has been
emitted, not at the moment mid-solution where the model's evidence goes
wrong.

The gap between the two lines is precisely a reward that pays an individual
\emph{in-reasoning} metacognitive act for the evidence it moves. \pmishift{}
supplies this: for each \metaopen{}\dots\metaclose{} block we measure the
model's gold-versus-decoy log-likelihood contrast immediately before and
immediately after the block under a frozen reference, and route the
resulting reward onto the block's tokens and nothing else
(Section~\ref{sec:method}). Knowing uncertainty is thereby rewarded chiefly
insofar as it is used: a block that merely announces doubt earns no
\pmishift{} credit, while one that redirects the solution toward the correct
answer earns the shift it caused. Our training objective does retain a small
explicit calibration term on the confidence token
(Section~\ref{sec:method:pmishift}), so we do not claim calibration arises as
a pure byproduct; rather, we treat accuracy as the primary outcome and
calibration (ECE, Brier, overconfidence rate; RQ4) as a secondary one---the
reverse of the calibration literature's emphasis, and, we argue, the right
ordering for reasoning.
